\documentclass{beamer}
\usepackage[utf8]{inputenc}
\usepackage[vietnamese]{babel}
\usepackage{graphicx}
\usepackage{tikz}

%--- Style riêng ---
\usetheme{Madrid}
\usecolortheme{seahorse} % xanh nhẹ
\usefonttheme{professionalfonts}
\setbeamertemplate{navigation symbols}{}

\begin{document}

%--- Trang bìa ---
\begin{frame}[plain]
  \centering
  \renewcommand{\baselinestretch}{1.1}\selectfont
  \vspace*{2cm}
  {\color{blue!80!black}\Huge\textbf{Giới thiệu Lý thuyết Đồ thị}}\\[0.5em]
  {\Large\textbf{Lê Văn Hoàng An}}\\[0.5em]
  {\today}
\end{frame}

%--- Slide 2: Giới thiệu ---
\begin{frame}{Giới thiệu}
  \begin{columns}[T]
    \begin{column}{0.6\textwidth}
      \begin{itemize}
        \item Lý thuyết đồ thị là một nhánh của toán học rời rạc.
        \item Nghiên cứu các đối tượng (đỉnh) và mối quan hệ (cạnh).
        \item Ứng dụng rộng rãi trong:
          \begin{itemize}
            \item Khoa học máy tính (mạng, thuật toán).
            \item Giao thông (đường đi ngắn nhất).
            \item Mạng xã hội (phân tích kết nối).
          \end{itemize}
        \item Đặt nền tảng cho nhiều thuật toán quan trọng.
      \end{itemize}
    \end{column}
    \begin{column}{0.4\textwidth}
      \centering
      \includegraphics[width=\textwidth]{Gemini_Generated_Image_dv0x2bdv0x2bdv0x.png}
    \end{column}
  \end{columns}
\end{frame}
